\chapter{Réalisation de Secure Requirements Studio}

\section{Environnement de développement}
\begin{itemize}
    \item Backend~: Python 3.12, FastAPI, Poetry pour la gestion des dépendances.
    \item Frontend~: Node.js, Next.js 15, TypeScript.
    \item Éditeur~: \TODO{VS Code / autre, selon l'environnement réellement utilisé}.
    \item Gestion de version~: Git. \TODO{Préciser l'hébergement du dépôt (interne,
    GitHub, GitLab...) et la politique de branches si applicable.}
\end{itemize}

\section{Architecture de Secure Requirements Studio}
Le backend est organisé en couches~: \texttt{routers} (endpoints HTTP), \texttt{services}
(logique métier des différents moteurs), \texttt{repositories} (accès aux données),
\texttt{models} (entités SQLAlchemy) et \texttt{domain} (connaissances métier statiques,
telles que le catalogue de concepts et les référentiels de sécurité).

\section{Questionnaire guidé}
Le \emph{Classic Questionnaire} permet de mener un entretien complet sans dépendre d'un
fournisseur d'IA~: chaque question est associée à un domaine du Knowledge Graph, avec
suivi automatique de la complétude et des dépendances entre concepts (par exemple, une
question sur le RBAC régional n'est posée qu'après confirmation du regroupement des
coopératives par région).

\section{Entretien assisté par IA}
En mode assisté, l'IA joue le rôle d'un consultant senior en ingénierie des exigences~:
questions courtes et précises en français, pas de répétition d'une question déjà posée
sous la même forme, et restriction stricte au périmètre du projet cadré. Toute question
hors périmètre reçoit une réponse de recentrage plutôt qu'une réponse de chatbot
généraliste.

\section{Intégration des fournisseurs IA}
\TODO{Décrire précisément, au moment de la rédaction finale, les fournisseurs IA
effectivement intégrés et testés (Gemini confirmé dans le code ; les autres
fournisseurs mentionnés dans les spécifications initiales — OpenAI, Anthropic Claude,
Kimi, Ollama — ne doivent être décrits comme livrés que s'ils ont réellement été
implémentés et testés au moment de la rédaction).}

\section{Génération dynamique des exigences}
Chaque réponse suffisamment complète déclenche la génération d'exigences typées
(fonctionnelle, non fonctionnelle, sécurité, technique, etc.), chacune associée à une
clé unique, un identifiant du concept source (\texttt{source\_concept\_key}), une
priorité et un statut — garantissant la traçabilité entre une exigence et la réponse
qui l'a produite.

\section{Génération du Cahier des Charges}
Le Cahier des Charges est assemblé automatiquement à partir de l'état courant du
Knowledge Graph et des exigences générées, avec pour contrainte de rester concis
(3 à 5~pages) plutôt que de reproduire un document générique volumineux.

\section{Génération de la conception}
La conception applicative de haut niveau (architecture, acteurs, modules,
flux principaux) est dérivée du même modèle de données que le Cahier des Charges,
garantissant la cohérence entre les deux documents.

\section{Génération du diagramme MVP}
Voir section~4.7 pour le détail du principe de génération du diagramme unique de
conception MVP.

\section{Génération des documents}
Les formats d'export pris en charge sont~: JSON (données machine), Markdown (source de
vérité textuelle), HTML (dérivé du Markdown), DOCX et PDF (générés directement à partir
du modèle de contenu structuré, pour garantir une mise en forme fiable indépendamment de
la conversion Markdown).
